\documentclass[10pt,a4paper,twoside,twocolumn,openany,bg=print,nodeprecatedcode]{dndbook}
\usepackage[utf8]{inputenc}
\usepackage{tabulary}
\usepackage[normalem]{ulem}
\usepackage{color,soul}
\usepackage{float}
\usepackage{caption}
\usepackage{wrapfig}
\usepackage{tocloft}
\usepackage[hidelinks]{hyperref}
\raggedbottom
\extrafloats{2000}
\cftsetindents{section}{1em}{0em}
\cftsetindents{subsection}{2em}{0em}
\cftsetindents{subsubsection}{3em}{0em}
\setcounter{tocdepth}{1}
\begin{document}
\addcontentsline{toc}{chapter}{Contents}
\tableofcontents
\newpage
\chapter{Diseases}
The Dungeon Master's Guide provides three example disease to use in a D\&D game which are simple to implement and easy to adapt to suit one's individual needs, but lack complexity or versatility. The diseases described below delve deeper into what a fantasy disease can be, and how it can be implemented into a story or campaign.
While the Medicine skill is always useful when dealing with diseases, some focus has been given to alternative skills to use when diagnosing or treating these diseases. Typically in D\&D a disease can be cured easily with magic at relatively low levels. In the examples below, the methods of treated are more difficulty or require additional steps to make the usual remedies effective, not only to make these conditions more dangerous to parties with easy access to disease-curing magic, but also to evoke the tense and dramatic process that medicine can often be in fiction.
In addition, some of these entries serve multiple purposes rather than simply being negative stat penalties to apply to players. For example, applying Blightpus to a horde of weak enemies can make them more dangerous to cleave through in close quarters, while Shroud Vine confers both penalties and desirable boons to those infected. At the end of this section are roll tables of plot hooks for each listed disease to use as inspiration for stories involving them.
 Each disease subheading is divided into the following parts:
 First the entry gives the most common name for the disease, but some disease may be known by multiple names or have a different name in other settings.
 Symptoms describes the visible effects of the disease to an outside observer as well as what an infected creature might feel or be aware of.
 Effects details the effects of the disease in game terms and how it affects an infected creature's in-game statistics.
 Contracting outlines how a creature becomes infected in game terms, as well as an overview of situations where one is likely to come into contact with the disease.
 Diagnosis gives suggestions on ability checks one could use to identify the disease or gain more information about it.
 Treatment describes how to cure the disease or how to alleviate some of its effects in game terms.

\begin{itemize}
\item Blightpus
\item Boiling Blood
\item Cannibal Craving
\item Contagious Troll Tumours
\item Earwyrms
\item Glass Transmorgification
\item Glass Delusion
\item Shroud Vine
\end{itemize}

\begin{DndComment}{Where are all the healers?}
In the magical worlds of D\&D, healing disease is simple and can be done with a single spell. Many large towns and cities have divine spellcasters available, and smaller settlements can call for aid. Wilderness and tribal communities may have druids in their ranks who are also competent healers. With healing magic everywhere, how does an outbreak of disease even happen? How can you run exciting and horrifying stories of plague-ridden lands when by all rights this should be a trivial matter to resolve?
The solution used for many of the diseases presented in this document is to simply make some of the diseases themselves more difficult to heal. But this may not work for your story. Perhaps you want the disease to be an integral part of the setting, but not something the players need to worry about for long, since the real-life countermeasures against disease spread such as staying at home aren't condusive to adventure. Or perhaps you simply don't want to have the hassle of repeatedly running complex cures that will slow down the pace of the game. In situations like this, the formal rules and their streamlined cures might work for you, but creates a problem when trying to explain why everyone can't do this. There are other possible explanations though:
The affected settlement is isolated and healers can't reach it, or the locals can't easily leave to call for aid.

\begin{itemize}
\item There are healers available, but not enough. They are being overwhelmed by the patient load, or the infection is simply too widespread to stop.
\item Political, religious, or racial divides means that certain demographics aren't being treated or are at a lower priority for recieving treatment.
\item Healers are split between treating this outbreak and providing aid for some other event, such as a war or natural disaster.
\item Healers are handling individual cases but can't address the underlying cause of the outbreak, such as contaminated water, biological warfare, attacks by infected monsters, or a magical curse.
\item Supplies needed to properly equip doctors and spellcasters are scarce.
\item A powerful authority is ignoring or blocking aid to the infected because doing so would detract from some political goal.
\item What if the players have one or more healers in their ranks, and wish to spend days performing healing services? It may be prudent to simply allow them to do this. Players have limited resources, and spending time healing means they have fewer resources available to confront danger, so this course of action comes with risk.
\end{itemize}

\end{DndComment}

\section{Plot Hook Tables}
Below are roll tables with a few suggested plot hooks or story ideas that incorporate the diseases described in this document. These are only suggestions and don't represent the whole scope of how a disease can be used in a campaign.
Glass Delusion and Glass Transmogrification use the same table.
\begin{DndTable}[header=Blightpus]{XX}

 & \\
1 & Miners are attacked by monsters infected with Blightpus, suggesting that their tunnels have crossed with caves that reach the underdark.\\
2 & An outbreak of Blightpus in a large town leads to some areas being quarantined, with the quarantined areas becoming rife with crime and adventurers are denied access to important people or places in the quarantine zones.\\
3 & Cultists to gods of disease or demons intentionally infect themselves with Blightpus to burn their foes in battle. \\
4 & A sharp surge in Blightpuff growth in nearby wilderness leads locals to fear that a powerful monster such as a dragon has claimed the area as their territory.  \\
\end{DndTable}


\begin{DndTable}[header=Boiling Blood]{XX}

 & \\
1 & Residents of an infected settlement seek the aid of a powerful mage who can control the weather, the cold temperatures allowing them to treat people en-masse.\\
2 & A healer accidentally kills a Boiling Blood victim via hypothermia while attempting treatment, and tries to cover it up.\\
3 & A red dragon has laired near a town, and infected water runs down from its mountain. Townsfolk want the dragon either slain or cured to halt the source of the infections.\\
4 & The scent of cooked flesh coming from the mass graves of Boiling Blood victims during an outbreak has attracted monsters to come and dig up the bodies.\\
\end{DndTable}


\begin{DndTable}[header=Cannibal Cravings]{XX}

 & \\
1 & A  suspected contamination of Cannibal Craving in an expedition's water supply causes paranoia and violent conflict.\\
2 & A cult of demon worshippers see the infection as a blessing from their deity and seek to spread it.\\
3 & A noble with Cannibal Craving murders and eats poor people in secret, covertly serving human meat to guests.\\
4 & A village is attacked nightly by frenzied animals infected with Cannibal Craving and seek aid in finding the source of the infection.\\
\end{DndTable}


\begin{DndTable}[header=Contagious Troll Tumours]{XX}

 & \\
1 & An adventurer returns to town after a successful quest and starts boasting about their newfound healing powers in a local tavern, unaware of the risk.\\
2 & Another monster contracted CTT and became more aggressive and confident from their enhanced healing.\\
3 & A renowned hero owes their success to CTT but hides their face from public viewing as they are covered in tumours.\\
4 & A troll ends up in an unlikely place (such as on a boat, in a locked prison cell, or a lord's manor) as a result of being born from a CTT victim.\\
\end{DndTable}


\begin{DndTable}[header=Earwyrms]{XX}

 & \\
1 & An Earwyrm repeats a dangerous secret it heard while occupying a previous host, launching a search to find the source of the outbreak.\\
2 & A stock of domesticated Earwyrms escapes, starting an infestation. Annoyed locals want the situation resolved by any means.\\
3 & Players need to seek out an Earwyrm breeder so they can use the Earwyrms to their advantage against a Mind Flayer colony.\\
4 & An Earwyrm host mistakes the creature's random mimicry for a cryptic supernatural portent.\\
\end{DndTable}


\begin{DndTable}[header=Glass Delusion/Transmorgification]{XX}

 & \\
1 & A noble is secretly suffering from Glass Delusion, and is being blackmailed by someone who knows their secret.\\
2 & Multiple cases of Glass Delusion in a single town pushes local authorities to investigate the possibility of a new potion black market\\
3 & A hag curses the loved one of an adventurer with Glass Delusion in revenge for a slight against her.\\
4 & A spellcaster's aid is needed for a quest, but they've become reclusive and paranoid due to suffering from Glass Delusion.\\
\end{DndTable}


\begin{DndTable}[header=Shroud Vine]{XX}

 & \\
1 & Shroud Vine bushes appearing near a countryside village are seen as a bad omen, and the villagers hire adventurers to investigate the source.\\
2 & An extremist druid intentionally infects people with Shroud Vine as punishment for perceived crimes until they can repay their wrongdoings by going on a quest.\\
3 & A dungeon is guarded by swarms of plant creatures that can only be avoided by infecting oneself with Shroud Vine.\\
4 & A settlement afflicted by a Shroud Vine outbreak is at odds with itself as conflict arises between those who want to cure the parasite and those who wish to coexist with it.\\
\end{DndTable}


\chapter{Curses}
This section details curses that can be used against players using the format for curses described in Van Richten's Guide to Ravenloft, where curses are generally described in three parts:
Pronouncement describes a threat or warning to the target of the curse, or anyone who might eventually fall foul to it.
Burden explains the effects of the curse in game terms.
Resolution explains how the curse might be ended.
\section{Bubblespeak}
This curse is often associated with fey for it's amusing nature. The affected creature dispenses bubbles from their mouth whenever they try to speak.
Pronouncement: This curse is levied at a creature that uses foul language in a highly inappropriate setting or who speaks ill of an important figure, and includes a reference to the old saying of 'washing your mouth out with soap'.
Burden: Whenever the cursed creature speaks, bubbles emerge from its mouth instead of words. If these bubbles are popped, the intended spoken words of the creature is emitted from it. Each bubble can contain a phrase up to 20 words long, so longer sentences create more bubble. The bubbles float upwards and are carried by even the softest breezes. Trying to pop a bubble requires a DC 13 Dexterity check.
Resolution: The curse instantly ends if the cursed creature uses soap to wash out the inside of its mouth.
\section{Cold of Heart}
This curse punishes callousness by trapping its victim in cold and loneliness with a literal heart of ice.
Pronouncement: Directed at someone who shows a lack of empathy. 'For your cold heart, may you never feel warmth again.'
Burden: The affected creature becomes unable to tolerate heat. If the creature is in an area with an ambient temperature higher than 0 degrees Celsius, it suffers the effects of Extreme Heat, but is immune to the effects of Extreme Cold. It gains a vulernability to fire damage, but is resistant to cold damage. Close contact with warm-blooded creatures, such as holding hands, grappling, or applying bandages, deals 1 fire damage to the cursed creature. The cursed creature itself has no body heat, and feels cold to the touch.
Resolution: Only an act of true kindness can break this curse.
\section{Collector}
A torment meant to punish greed, or to reduce a powerful foe into an ineffective eccentric. The cursed is driven to collect a certain kind of object.
Pronouncement: The pronouncement of this curse references a specific kind of object. It could be part of a set such as every shard of a broken sword or every coin from a stolen treasure, or it could be a kind of general object such as gems, spoons, mugs, dice, figurines, shoes, etc.
Burden: The affected creature feels a compulsion to collect the specific type of object. It refuses to relinquish objects of that type in its possession or allow objects in its posession to be destroyed. If another creature takes an object from the cursed creature's collection, it is magically compelled to retrieve the object by any means available to it. If the creature goes a day without gaining a new object for its collection, it must make a DC 10 Wisdom saving throw. On a failed save, it gains one form of indefinite madness (DMG page 260). For each day the creature goes without aquiring a new object for its collection, or retrieving an object taken from its collection, the DC increases by 1. After it gains a new object, the creature loses one form of madness every 24 hours and the saving throw against madness decreases by 1, as long as it continues to grow its collection.
Resolution: If the objects belong to a set, then returning the completed set to its rightful owner will lift the curse. Knowing this doesn't affect the cursed creature's behaviour: it still refuses to part with its collection.
For a generic object type, a creature can be freed from this curse if it destroys its own collection. Knowing this doesn't change its behaviour, it still refuses to part with its collection in any way, but it can still be freed if it unknowingly or accidentally destroys its collection, or if it is compelled to do so by magic. The creature has to destroy its own collection: another creature doing so won't relieve the burden of the curse.
\section{Duplicate}
This devious curse creates an exact duplicate of its target that can torment it in a variety of ways, from being a perfect match in battle to more insidious psychological attacks.
Pronouncement: This curse affects a creature that disturbs an ancient sarcophagus containing the soul of an ancient sorcerer, imprisoned fiend, or other villain. An inscription on the tomb warns against tampering, lest 'You shall be your own worst enemy'.
Burden: A creature affectd by this curse suffers no immediate effects. 24 hours after the curse takes effect, an exact duplicate of the cursed creature appears inside the sarcophagus. It has all the memories of the cursed, along with all of its game statistics and class features, but none of its equipment. However, it despises the cursed creature, and either wishes to ruin that creature's life and reputation, or wishes to kill and replace the cursed and live on in their place. If the duplicate is killed, another exact duplicate appears in the sarcophagus with all the memories of the previous duplicate.
Resolution: Killing the duplicate is only a temporary measure. The only way to end the curse is to re-seal the duplicate inside the sarcophagus.
\section{Eternally Forgotten}
As a punishment against the vainglorious, people cursed in this way can never achieve lasting recognition for their deeds, and even simply interpersonal relationships become impossible.
Pronouncement: This curse is directed at those who seek to be famous and commit terrible deeds to do so. Alternatively, it is directed at those who seek to rewrite history in their favour.
Burden: If the cursed creature isn't percieved by someone else for 10 consecutive minutues, that person forgets that the cursed person ever existed. Any illogical or impossible memories as a result of this are rationalised e.g. if the cursed creature did something significant, other people simply assume it didn't happen or that someone else did it. Other people tend to ignore or overlook evidence of the cursed person's existence unless it is explicticly pointed out to them.
Resolution: The curse is lifted if the cursed intentionally takes a course of action where they turn down the possibility of renown of their deeds, or allow someone else to take credit for their deeds. For example, they could destroy a statue erected in their honour, or claim that a companion actually slew a monster they killed.
\section{Horrid Reflection}
This curse is only noticable in the reflection of the afflicted, revealing their true nature.
Pronouncement: This curse is levied at someone with an evil nature whose wicked deeds are kept secret, usually by someone who knows the secret but cannot expose it, or whose testimony won't be believed.
Burden: The affected creature experiences a nightmare where they see themselves twisted in some monstrous way, such as appearing like a corpse, or having bestial features. When the creature awakens, it appears normal. However, it's reflection always appears as it did in the nightmare, to anyone who sees it. If the cursed creature sees its reflection, it becomes frightened of it. If it attacks its own reflection, such as by shattering a mirror or throwing a stone into water, the cursed creature takes 5 (1d10) psychic damage.
Resolution: Publically revealing the terrible secret or true nature of the afflicted does not end this curse. Instead, it causes the cursed creature's appearance to become like that of its reflection. Only true atonement, and a sincere change of heart, will end this curse. Once the creature's nature has changed, it ceases to be frightened by its reflection, and touching its reflection returns it to normal.
Alternatively, if the cursed creature reveals its own villainy, and fully embraces its deeds, then the curse's magic fails: the creature's wickedness is no longer a secret burden, thus this curse is no longer able to serve as a suitable punishment.
\section{Nonsense}
A tongue-twisting curse that inhibits a creature's ability to speak.
Pronouncement: This curse is caused by reading a cursed book written in bizzarre nonsensical script, and the words 'The words of this book make no sense. Read them, and only Nonsense will make sense. So be Sensible'.
The inside cover of the book warns against attempting to decipher the scrawl within. Merely observing the pages isn't enough, a creature only becomes cursed if they attempt to translate, decode, or otherwise interpret the words written on the pages.
Burden: The cursed creature cannot communicate with or understand any language. If it tries to speak or write, it only knows how to communicate in Nonsense. Translating nonsense is impossible: it has no translation, since it is completely nonsensical. Telepathy does not allow sensical communication either. However, two people cursed in this way can understand each other perfectly through any medium, even if they would normally not share any languages.
Resolution: A creature affected by this curse can still learn new languages, since languages it previously didn't understand are nonsensical to them and thus aren't affected by the curse's burden. Doing so creates sense out of nonsense, breaking the curse's magic. The same applies if it uses magic to be able to understand a language it previously couldn't (but not if another creature attempts to use magic to translate Nonsense, since that doesn't work, or if it uses magic to understand a language it is normally proficient in). Effects that completely disable a creature's ability to understand any kind of language (such as the Feeblemind spell) cause the curse to fail.
\section{Rat Lord}
Rats follow in the footsteps of the cursed, whether they like it or not. As rats bear diseases and pilfer food, this is likely to make the cursed unwelcome in humanoid settlements.
Pronouncement: This curse is usually directed at people of high social status or noble title, such as a knight or respected priest, and is done so in response to the individual committing some dishonourable deed not in line with the expectations of their station, with the utterance of 'If you must act so wretchedly, then you deserve only to be a master of rats.'
Burden: Rats within 300ft of the cursed creature become obsessed with them, and belive the cursed to be their Lord. These rats always know the direction to the Rat Lord and do their best to follow them or get near to them. Once a rat has been influenced in this way, the effect persists until the curse ends or the rat is more than 1 mile away from the Rat Lord for 24 hours. The rats will flee if attacked, but return as soon as they believe the danger to be over, and will only flee as far as they need to. Otherwise, their only goal is to be as close to their lord as possible at all times.
Resolution: The cursed must complete an act of nobility in line with their noble rank or social status, usually involving some sort of sacrifice or an act of great humility to prove themselves worthy of their station.
Alternatively, they can pass on the curse to someone else. A willing creature declares themselves the new Rat Lord, while the cursed makes a symbolic offering and humbles themself before the new lord.
\section{Two-Faced}
This curse is meant as an ironic punishment - and a horrifying transformation - towards traitors and deceivers.
Pronouncement: A person that suffers a terrible fate as result of a betrayal or lie told about them may lay this curse through their anguish and rage.
Burden: The cursed creature grows a second face on their body, usually out the back of their head. Whenever a creature tells a deliberate lie or tries to conceal the truth, the second face will reveal the truth. If the cursed creature tries to persuade someone of the truth or reveal a lie, then the second face will tell a lie.
Resolution: This curse ends when the cursed creature reveals the truth about the lie they told or proves the innocence of someone they betrayed. They must also reveal the presence of their curse, laying bare their two-faced nature.
\section{Vanity's Reward}
This curse is meant as an ironic punishment for those who overly value physical beauty.
Pronouncement: Usually directed at someone who has just committed a grevious act of vanity, or disparages someone for their superficial appearance, especially if that person is much mightier than their looks would suggest. If beauty be all that concerns you, then let it be eternal, cast in stone!'.
Burden: The cursed creature slowly begins to turn to stone. At first only visual indicators appear, with patches of stony-looking skin appearing on their body. After one day their speed is halved and they have disadvantage on all attack rolls, ability checks, and saving throws that use Dexterity. After two days their speed becomes 0 ft and can't benefit from increases to speed. After three days, they become Petrified.
Resolution: Magic that would cure petrification can't end this curse. Instead, it can only revert the creature to a partially cured state, as if it had the curse for one day. To permenantly end the curse, the creature must either intentionally deface its physical appearance, or intentionally appear in public in a humiliating or unappealing state.
\section{Wayward Shadow}
This curse is usually levied at those who committed a terrible crime and got away with it. The shadow of their deeds shall follow them forever.
Pronouncement: This curse is uttered by the victim of a crime who knows that the crime against them will never be solved, such as a dying person who knows there are no witnesses, or a marginalised beggar who knows that nobody will bother to seek justice on their behalf. 'Your wickedness shall forever follow in your shadow, but remember: what happens in the dark always comes to the light.'
Burden: The cursed creature's shadow becomes magically animated. Once a day, it can leave the cursed creature's side, and uses the statistics of a Shadow (Monster Manual pg 269). If it is destroyed or is away from the cursed creature for more than 24 hours, it reappears at the creature's side, as if it were an ordinary shadow. While independent, it seeks to commit evil acts similar to that committed by the cursed creature. It may also attempt to implicate the cursed creature in these crimes. At all other times, it appears to be an ordinary shadow, and cannot be detected by magic that would detect undead. While away, the cursed creature casts no shadow.
Resolution: The curse can lifted by applying the words of the pronouncement in a literal or figurative manner. Solving the original crime and bringing that knowledge to light will lift the curse, but only if the wrongdoer would suffer some kind of punishment for doing so. Applying the words literally can also end the curse: catching the wayward shadow in the midst of its misdeeds and casting bright light over it destroys the Shadow, and ends the curse.
The Atonement function of the Ceremony spell can also end the curse, if it is cast by a follower of a deity of light or the sun.
\begin{DndComment}{Non-player curses}
What happens when the target of a curse isn't a player? A curse can be a central element in a story. Perhaps the players are tasked with lifting a curse, or the villain of the story is the curse's victim. When not applied to players, you can take greater liberties with how the curse affects its victim. As written, a Rat Lord has no control over the rats that follow them, but you may wish a villainous Rat Lord to have such control. Or a non-player cursed with a Horrid Reflection frightens other people instead of themselves with their reflection. Such abilities could be too strong in the hands of creative players, and thus the curse fail to be a meaningful burden at all. But in the hands of the DM, can enhance the story. Many a curse in fiction comes with great power attached, but usually in such a way that the power isn't worth the price.
You can also use this to make a curse much more of a burden on a non-player character. Some of the mechanics of these curses as listed are meant to be fair on players, still allowing a player to enjoy the game while cursed, or to mitigate some of the curse's worst effects. NPCs under the control of the DM need no such mercy, and thus their curse might be more severe or more difficult to resolve.
\end{DndComment}

\chapter{Dungeon \& Wilderness Hazards}
\section{Flora \& Fauna Hazards}
This kind of hazard includes all kinds of dangerous plant, fungus, or immobile beasts.
Bloomers
Bursting Blightpuff
Cave Mist Creeper
Coffin Moss
Giant Clams
Pixie Sniffles
Pond Snare
Scylla's Claw
\section{Natural Hazards}
This kind of hazard or terrain abnormality come about as a result of geological phenomena and typically don't involve magic. 
Brine Pools
Lodestones
Magic Lodestones
Natural Gas
\section{Magical Phenomena}
These are supernatural events, and only arise from the presence of magical beings or incursions from other planes.
Manes Vapour
Stain of the Far Realms
\chapter{Traps}
Ballista Trap
Burning Oil Trap
Electric Water
Fish Hook String
Gong
Greased Stair
Psychic Disruptor
Razor Blades
Razor Tripwire
Rigged Lantern
Shoving Statue
Slashed Bridge
Spike Plates
Spring Loaded Door
Wind Tunnel
\begin{DndComment}{Glyphs of warding as traps}
An easy way of creating a magical trap is to use the Glyph of Warding with the Spell Glyph function. This can be an actual Glyph of Warding spell, or other magic that functions in an identical manner, and makes sense in any dungeon that is guarded by a villain with access to spells of 3rd level or higher. But what spells to use? A simple damaging spell such as Fireball or Cloudkill is the most obvious choice, but there are endless creative options that let you use this kind of trap in interesting ways. Consider some of the following spells:

\begin{itemize}
\item Entangle near enemies or other traps.
\item Grease at the top of a steep slope.
\item Thunderwave placed so enemies can hear the spell effect.
\item Enlarge/Reduce so that a creature can't fit out the door of the room it's in.
\item Gust of Wind near a cliff or trap.
\item Silence set so that it only triggers when someone attempts to cast a spell.
\item Polymorph in a situation that requires hands or speech.
\item Antilife Shell to block a passage guarded by constructs or undead.
\item Contagion near the beginning of a dungeon to weak intruders in advance.
\item Danse Macabre cast on guards that die near the glyph.
\item Modify Memory combined with Magic Mouth in a passage to create a false memory of finding a dead end.
\item Passwall to create pitfalls in solid floors or allow enemies to enter a room.
\item Flesh to Stone to trap intruders.
\item Reverse gravity to place pit traps on the ceiling, or cause them to trigger twice.
\item Earthquake to destroy a trapped structure with intruders inside.
\end{itemize}

\end{DndComment}

\chapter{Siege Equipment}
The Dungeon Master's Guide describes seven kinds of siege equipment, and mentions how some kinds of siege weapons can be loaded with different types of projectiles for different effects. This section describes five more kinds of siege equipment based on real historical or mythical devices, as well as a number of examples for alternative projectile types, including ones mentioned but not described in the DMG.
\section{Additional Siege Equipment}
Burning Bale
Burning Glass
Flamethrower
Hwacha
Mobile Barricade
\section{Additional Projectile Types}
Ballista
The large bolts of a ballista are easily attached to ropes or chains and fitted with barbed heads to create harpoons ideal for attaching to ships and larger creatures such as giants or whales, either to allow creatures to climb along the ropes or to pull the target closer.
In settings with gunpowder technology, ballista bolts can be fitted with fireworks that increase the range of the weapon and detonate shorty after being fired.
Ballista (Firework)
Ballista (Harpoon)
Cannon
The cast iron balls of a cannon are devestating testaments to the power of gunpowder, but technology also allows for more flexibility beyond these simple blunt projectiles.
Bundles of smaller projectiles, known as Grapeshot, are effective at close-range obliteration of crowds, while also being easier to improvise. An ill-equipped force can stuff a cannon with cutlery and pebbles when they run out of proper ammunition.
Two cannonballs connected by a chain whirl through the air and wrap around their targets, designed for shredding the sails and rigging of ships, they also work well against flying foes.
Fitting the propellant of the cannon into the projectile rather than into the barrel allows it to fire explosive  rockets with a shrapnel payload.
Cannon (Chain Shot)
Cannon (Grapeshot)
Cannon (Rocket)
Suspended Cauldron
Acid is an effective alternative to boiling oil, especially against foes resistant to fire or when there are no facilities to heat the oil.
Grease isn't directly harmful to creatures it hits, but it leaves a slippery puddle behind that impedes anyone trying to charge past the defenders as they re-arm the cauldron, as well as being highly effective against foes trying to breach a fortified position by climbing ladders.
Icy cold water is very easy to source in certain climates, making it an economical choice of weapon. Against foes attacking in cold windy weather, it can quickly cause hypothermia, and leave a slippery hazard as it freezes on surfaces.
Though unlikely to cause serious harm, sewage is best used as a deterrent to demoralise a foe or to cause nausea that reduces the effectiveness of their assault. Often used as a desperation measure by besieged foes that have expended their more dangerous fluids.
Suspended Cauldron (Acid)
Suspended Cauldron (Grease)
Suspended Cauldron (Ice Water)
Suspended Cauldron (Sewage)
The Mangonel and Trebuchet both hurl the same kinds of projectiles via different methods, but their range and damage varies.
Hay soaked with oil can be set alight prior to firing. Igniting the hay takes place during the action used to load the weapon. Less effective at destroying fortifications than a simple stone except those made of wood.
Catapults make for an excellent delivery method for primative biological warfare. Baskets or hives full of stinging insects can quickly diminish enemy morale and cause disarray among the defenders, while festering corpses of humanoids or animals are a potent weapon of intimidation as well as a vector for disease that is especially effective during drawn-out sieges. Normally these diseases are represented by Sewer Plague (described on page 257 of the DMG) but can be almost any sort of disease, including ones found in this document such as Boiling Blood or Blightpus.
Wheels or barrels embedded with blades can roll for a time after being hurled, they lack the impact of a boulder but are much better at slicing through tightly-packed infantry.
Mangonel
Mangonel (Flaming Ball)
Mangonel (Insect Hive)
Mangonel (Corpses)
Mangonel (Blade Wheel)
Trebuchet
Trebuchet (Flaming Ball)
Trebuchet (Insect Hive)
Trebuchet (Corpses)
Trebuchet (Blade Wheel)
\chapter{Poisons}
The poisons presented here are an addition to those presented in the Dungeon Master's Guide. 
\begin{DndTable}[header=Poisons]{XXX}

Item & Type & Price per dose\\
Blackout Brew & ingested & 400 gp\\
Carrot Soup Special & inhaled & 150 gp\\
Concentrated Capcaisin & ingested & 50 gp\\
Cyanide Pill & ingested & 750 gp\\
Death's Slumber & ingested & 50 gp\\
Ghost's Touch & injury & 200 gp\\
Itching Powder & contact & 100 gp\\
Laughing Gas & inhaled & 200 gp\\
Psilocybin & ingested & 400 gp\\
Quickbleed & injury & 1200 gp\\
Stingeye & contact & 150 gp\\
Ten Step Death & contact & 1200 gp\\
Theif's Red Hand & contact & 300 gp\\
\end{DndTable}


Blackout Brew
Carrot Soup Special
Concentrated Capsaicin
Cyanide Pill
Death's Slumber
Ghost's Touch
Itching Powder
Laughing Gas
Psilocybin
Quickbleed
Stingeye
Ten Step Death
Theif's Red Hand
\chapter{Bestiary}
There are countless ways of applying the contents of this document in your adventures. One way is to use them to inspire new monsters or as modifications to existing monsters. Sometimes a small twist is all you need to make a familiar foe feel fresh and exciting. Below are a few sample stat blocks to aid you. They are all creatures from the Monster Manual with one or two appropriate alterations. Bear in mind a couple of thing: First, that these are just examples to get you started and don't even begin to scratch the surface of how curses, diseases, hazards, traps, or poisons could be used in a campaign. Second, that non-player characters need not follow the same rules as player characters. Instead of following the rules written here to the letter, your focus should instead be on creating something that makes intuitive sense to the players, and is fun to engage with. For example, Contagious Troll Tumours only grants a very small amount of hitpoint regeneration to an infected player, since regeneration is a very powerful ability, but isn't meaningful when applied to a monster in a combat encounter. Instead, it makes more sense to give that monster regeneration more akin to that of a troll. It's more impactful and makes more sense to the players. Your players can't see the monster stat blocks, and likely won't notice as long as the mechanics remain somewhat consistent.
Awakened Shroud Vine Shrub
Shroud Vine bushes grow into eerie silhouettes of the creature that spawned them, but are otherwise normal bushes. Thus, they can be targeted by any magic that affects plants, such as the Awaken spell.
Shroud Vine Druid
While Shroud Vine is an insidous blight to most people, some Druids see it as a means of connecting with nature, and form a symbiotic relationship with the parasite, and wield it as a weapon against their enemies.
Blightpus Kobold
Kobolds dwelling in the underdark are at great risk of Blightpus outbreaks due to their cramped lairs and communal living. Infected kobolds are likely to charge recklessly at threats, sacrificing themselves to use their infection against enemies of the tribe.
Tumourous Cave Bear
When wildlife or monsters ingest troll meat, they risk contracting Contagious Troll Tumours. While their increased resilience is a temporary boon that may lead to them dominating the local ecosystem and hedging out rivals, the disease will eventually overcome them, leading to the birth of a new troll. 
Noble of the Horrid Reflection
Of the people likely to be cursed in this way, nobles are especially prime targets as they depend so heavily on their reputation, and the terror of having one's true nature revealed can eat away at them. Nonetheless, some villains aren't afraid of their wicked nature, but their foulness terrifies others, leaving them to fester in evil isolated from others.
Two-Faced Imp
As befitting their role as tempters and deceivers, devils are more likely than anyone to be cursed with two faces. Some come under this kind of curse as a result of clerics or divine entities that oppose fiends, other devils are cursed as punishment by their superiors.
Lodestone Elemental
Sometimes the forces that create  Lodestones can also give rise to elemental creatures. They are potent guardians of druidic temples, especially against foes that wield metal.
Shambling Mist Mound
Shambling mounds accumulate all kinds of plant matter as they devour their way through the wilderness. CMC normally stays static, feeding on the carrion its vapours kill, and ordinarily wouldn't fare well with mobility. However, a Shambling Mound generates enough carrion by itself that the fungus can easily survive.
'False Dragon' aka Flamethrower Roc
Rocs, like crows, can learn how to use tools, so this flamethrower device is activated by a simple switch that the Roc can operate itself. The False Dragon is an attempt to create a flying artillery weapon with the might of a dragon, capable of devestating fly-by strafing runs, but one that also won't grow up to be an intelligent titan that can resent its captivity. Only two kinds of people would create a False Dragon: people know know all there is to know about Rocs, and people who know nothing about Rocs.
Fetid Hill Giant
Hill Giants are notoriously filthy, and have no dignity in battle. If you think you've caught a giant off-guard in the latrine, you're in for a revolting lesson. A hill giant has no shame in tossing the contents of its chamber pot at attackers, then clubbing the unfortunate creature with the pot, or hurling it at foes outside of the splash zone.
Venomous Assassin
Poison is an incredibly potent tool for assassination. Some assassins master the art of using poisons, and may pay tribute to gods of poison and death in order to learn their arts. Wyvern venom is common for its high toxicity, but some assassins use other mixtures. Quickbleed is useful against foes that may try to flee or take cover, as they can bleed to death even after leaving the assassin's sights. Assassins also know that they could suffer terrible fates if captured alive so carry cyanide pills to ensure they cannot be tortured, often concealed in a false tooth. 
\end{document}
